
\titulo{9}

De niños me miraste dulcemente \\
y nos enamoramos: nos temblaban los músculos, \\
los ojos se nos volvían remansos \\
y no nos aguantábamos las ganas de abrazarnos como locos. \\
Pero la vida nos lanzó a piedrazos \\
y hacía veinte años que ya no nos veíamos las caras. \\
Pasábamos los días mirando compulsivamente pantallas, \\
mensajes codificados con luces que nos quitaban el sueño, \\
descripciones simbólicas del estado exacto del universo, \\
de calles empedradas con el rompecabezas de adoquines \\
y el mito urbano de la higuera en flor. \\
Y pese a que seguíamos creyendo \\
en ese mundo al que nos referíamos, \\
ya nunca transitábamos las largas avenidas, \\
los árboles frutales quizás estaban secos. \\
La realidad se había convertido \\
en una hipótesis innecesaria. \\
Navegábamos días de representaciones \\
que eran la verdadera y única realidad. \\
Y mirando los símbolos \\
que ya no significan más que símbolos \\
que ya no significan más que símbolos \\
me la paso esperando respuestas que no llegan: \\
que alguien prescindirá de mis servicios \\
y engañaré el estómago con unos mates tibios, \\
que hoy es tu velatorio y el entierro es mañana \\
y en todos estos años \\
no me animé a decirte que te amaba. \\
La poesía genuina no está ni en las pantallas ni en los libros, \\
ni en las recitaciones de poesía: \\
es el ``Raquel te amo'' \\
rayado con la birome sin tinta \\
en la puerta del inodoro público. \\

